% SNAS 2026 double-blind abstract draft.
% Compile with XeLaTeX so the document uses the installed Times New Roman font.
\documentclass[12pt,letterpaper,twoside]{article}

\usepackage[margin=1in]{geometry}
\usepackage{fontspec}
\setmainfont{Times New Roman}
\usepackage{setspace}
\doublespacing
\usepackage{fancyhdr}
\usepackage[hidelinks]{hyperref}
\hypersetup{
  pdftitle={Beyond Stable Answers: Abstract Draft},
  pdfauthor={},
  pdfsubject={Double-blind draft for the 5th SNAS Interdisciplinary Research Conference},
  pdfkeywords={explainable artificial intelligence, vision-language models, construction safety}
}

\setlength{\parindent}{0.5in}
\setlength{\parskip}{0pt}
\pagestyle{fancy}
\fancyhf{}
\fancyhead[LE,LO]{\textbf{DRAFT - NOT CAMERA READY}}
\fancyhead[RE,RO]{\thepage}
\renewcommand{\headrulewidth}{0.4pt}
\setlength{\headheight}{15pt}

\begin{document}

\begin{center}
\textbf{Beyond Stable Answers: Auditing Explanation Faithfulness in a Construction-Safety Vision-Language Pipeline}
\end{center}

\begin{center}
\textbf{Abstract}
\end{center}

Vision-language models (VLMs) can support construction-safety monitoring by connecting natural-language safety rules to visual evidence. However, answer accuracy alone does not establish whether the evidence underlying a safety judgment is faithful or robust. This feasibility study adapts a six-dimensional explainable artificial intelligence evaluation framework, originally developed for tabular intrusion detection, to a VLM-grounded construction-safety pipeline. We evaluated Florence-2-base-ft on 163 held-out images from ConstructionSite 10k representing compliant scenes, personal protective equipment violations, fall hazards, and struck-by risks. Because Florence-2 does not provide native free-form visual question answering in this implementation, open-vocabulary grounding identified workers and rule-relevant safety objects, while a deterministic person-relative rule produced the compliance judgment. Explanation quality was tested through rule-aware region ranking, safety-object masking, severity-controlled perturbations, targeted and random-location occlusion, size-invariant centroid drift, bounding-box disappearance, worker-loss correction, bootstrap confidence intervals, and paired Wilcoxon tests. Area ranking selected the worker rather than the safety object in 50.3\% of usable cases; rule-aware ranking reduced this to 0.6\%. The corrected rule-aware top-region ablation changed 37.4\% of judgments. Targeted object occlusion changed 39.2\% of judgments and produced greater explanation drift than the strongest random-location control (0.198 versus 0.142; paired \textit{n} = 144, \textit{p} = .0011). Although intersection over union initially classified 25.0\% of stable-answer cases as drifted, only 13.7\% moved more than 0.2 of the image diagonal, revealing substantial small-box bias; 5.1\% of targeted explanations disappeared entirely. These findings show that trustworthy construction-safety AI requires explanation-level evaluation, semantic region selection, size-aware metrics, and explicit failure accounting in addition to answer-level performance. The pipeline remains a research testbed requiring human oversight, not an autonomous safety system.

\end{document}
